% Part: proof-theory
% Chapter: sequent-calculus
% Section: rules-G3c

\documentclass[../../../include/open-logic-section]{subfiles}

\begin{document}

\begin{table}
\begin{tabular}{@{}c@{}c@{}}
\multicolumn{2}{@{}c@{}}{بديهي اصول}\\[2ex]
$!C, \Gamma \Sequent !C$ & $\lfalse, \Gamma \Sequent \Delta$
\\[2ex]
\multicolumn{2}{@{}c@{}}{منطقي قواعد}\\[2ex]
\Axiom$\lnot !A, \Gamma \fCenter !A $
\RightLabel{\LeftR{\lnot}}
\UnaryInf$\lnot !A, \Gamma \fCenter \Delta$
\DisplayProof
&
\Axiom$!A, \Gamma \fCenter$
\RightLabel{\RightR{\lnot}}
\UnaryInf$ \Gamma \fCenter \lnot !A $
\DisplayProof
\\[4ex]
\Axiom$ !A, !B, \Gamma \fCenter \Delta$
\RightLabel{\LeftR{\land}}
\UnaryInf$ !A \land !B, \Gamma \fCenter \Delta$
\DisplayProof
&
\Axiom$\Gamma \fCenter !A$
\Axiom$ \Gamma \fCenter !B$
\RightLabel{\RightR{\land}}
\BinaryInf$ \Gamma \fCenter !A \land !B $
\DisplayProof
\\[4ex]\Axiom$!A, \Gamma \fCenter \Delta$
\Axiom$!B, \Gamma \fCenter \Delta$
\RightLabel{\LeftR{\lor}}
\BinaryInf$!A \lor !B, \Gamma \fCenter \Delta$
\DisplayProof
&
\Axiom$\Gamma \fCenter !A$
\RightLabel{\RightR{\lor}}
\UnaryInf$ \Gamma \fCenter !A \lor !B$
\DisplayProof
\\[3ex]
&
\Axiom$\Gamma \fCenter !B$
\RightLabel{\RightR{\lor}}
\UnaryInf$ \Gamma \fCenter !A \lor !B$
\DisplayProof
\\[4ex]
\Axiom$!A \lif !B, \Gamma \fCenter !A$
\Axiom$!B, \Gamma \fCenter \Delta$
\RightLabel{\LeftR{\lif}}
\BinaryInf$!A \lif !B, \Gamma \fCenter \Delta$
\DisplayProof
&
\Axiom$!A, \Gamma \fCenter !B$
\RightLabel{\RightR{\lif}}
\UnaryInf$ \Gamma \fCenter !A \lif !B $
\DisplayProof
\\[4ex]
\Axiom$ !A(t), \lforall[x][!A(x)], \Gamma \fCenter \Delta$
\RightLabel{\LeftR{\lforall}}
\UnaryInf$ \lforall[x][!A(x)],\Gamma \fCenter \Delta$
\DisplayProof
&
\Axiom$ \Gamma \fCenter !A(a) $
\RightLabel{\RightR{\lforall}}
\UnaryInf$ \Gamma \fCenter \lforall[x][!A(x)]$
\DisplayProof
\\[4ex]
\Axiom$ !A(a), \Gamma \fCenter \Delta $
\RightLabel{\LeftR{\lexists}}
\UnaryInf$ \lexists[x][!A(x)], \Gamma \fCenter \Delta$
\DisplayProof
&
\Axiom$ \Gamma \fCenter !A(t) $
\RightLabel{\RightR{\lexists}}
\UnaryInf$ \Gamma \fCenter \lexists[x][!A(x)]$
\DisplayProof
\end{tabular}
\caption{د شهودي منطق د \Log{G3i} قواعد۔ په $\RightR{\forall}$ او $\LeftR{\exists}$ کښې !!{constant}~$a$ بايد په وروستي سېکوېنټ کښې رانۀ شي۔ $\Gamma$ د !!{formula}s څوګونې مجموعه ده؛ $\Delta$ تشه کېدے شي يا يو !!{formula} لرلے شي۔ د پېژندنې په بديهي اصولو کښې !!{formula}~$!C$ بايد اتومي وي۔ \Log{G1m} د \Log{G1i} هغه بڼه ده چې \RightR{\Weakening} او د $\lfalse \Sequent \quad$ بديهي اصل نۀ لري۔}
\ollabel{tab:G3c}
\end{table}

\end{document}